% Event Driven Systems
\cite{chavan2024fault}
As systems become larger the likelyhood of failure is stronger, in event driven systems specifcially, it becomes more significant with the asynchronous and distributed nature.
Reliability is important for these types of architectures as the ability to remain functional in the face of failure is essencial. Communication channels speciufificly are vulnerable to communicaition loss, latency
fault tolerance can be achieved through redundancy and replication, idempotency, circuit breakers (what I'm doing - observe functionality and suspend request to flawed services within a certain threshold of errors)


% Event Driven Architecture
\cite{manchana2021event}
emphasizes the production, detection, consumption, and reaction to events
event is a significant occirance or state change within a system that requires action
events are the primary means of communicaition between different components in the system. As a result the components can operate independantly and asynch
events - unit of information
architecture contains - event publishers, event mediators, event producers, event consumers, event brokers, event processing, event aggregation/correlation, and event schemas
Events allow for system intergration through exchange and synchronization between different systems, and events generated by sensors/devices can be processed in real time
When using events in modern systems then even schema should be well structured for interoperability and efficient routing mechanisms should be in place to ensure consumers only recieve what they are interested in. There should be error handling mechanisms to deal with message loss, duplication, and out of order delivery. There should be monitering and logging to have insight into the flow of events


\cite{kommera2020power}
paradigm enables systems to react to real time events by triggering actions in response to state changes
an event is a significant occurance in the system that can trigger actions or changes to components
architecture is built around asynchronous communication between components, so events drive system behaviour instead of relying on continious polling
event producers generate events and event consumers listen and respond to events
event handlers are responsible for processing events when they are detected (ex model transformations)
asynch communication allows components to operate without waiting for a direct response and decouples integration between components, allowing for horizontal scaling




% State Estimation in Event Driven Systems
\cite{martinez2015board}
event based state estimator - something is monitering a measured variable in order to detect an event. An estimator on the other side waits for events in the communication channel and uses them when they arrive
Authors have previously proposed estimators that work in the absense of information from sensors, triggering when theres no updates 
They propose a solution that triggers an estimator when the error is below a certain threshold
in ther case a consumer triggers a sensor to send data on a scheduel and when no recieved uses estimated values
state estimation itself can be event driven and not periodically sampled. their estimator does not assume continious or peroidict measurements and they arrive at irregular timestamp
estimation advancers through prediction and correction steps





% Fault Tolerance
\cite{hanmer2013patterns}
fault tolerance is being able to operate though some part is no longer operating correctly, so the book focuses on mechanisms that allow for continued operation, focused on improving the reliability and avaliability of software 
fault error failure chain
faults - defect present in the system that can cause an error. Represents a deviation from correctness. 
error - incorrect system behaviour, categories: value error - incorrect discrete values or system ststae, and timing - non performance time, caused by faults
failure - system behaviour is not performed as expected/according to system specifications, caused by errors
- grouped into consistent (appears the same each time it was observed), or inconsistent (appear different to different observers)
System Reliability - probability that it will perform without deviations from expected behaviour for a specific peroid of time, ie no failures during a specific timeframe
- measured with mean time to failure (MTTF) and mean time to repair (MTTR)
Dependabulity - measure of system trustworthyness to be rleied upon - reliability, avaliability, safety, recurity
Markov models can be used to show the transitions between failure and repair states
performance is linked closely to reliability
phases to fault tolerance - error detection, error recovery, error mitigation, and fault treatment
Error Recovery - undoing the bad effect and creating and error free state in the system so it can resume execution
Recovery Patterns
Quarantine=isolate element to prevent corruption of the rest of the system
Concetrated Recovery = remove distractions during the error
Error Handler = controlled mannor for handling errors
Restart = resume execution by restarting program from begining 
Rolleback = resume execution by moving to a state before the error occuranced
Roll Forward = resume execution by advancing to a future state that would have been reached if no error occured
Return to Reference Point = resume execution by returning to a specifci known state that is safe
Limit Retries = do not return to the error without changing something, as the error can reoccur
Failover = recovery by switching to a redundant data unit
Checkpoint = save state periodically so that it does not need to be regenerated
What to save = checkpoints save information of global interest for long duration processes
Remote Storage = depends on where you place checkpoints
Individuals decide timing = let each processes decide when to take a checkpoints based on their own needs
Data Reset = restore some data to an initial or predetermined value when it is incorrect


\cite{avizienis2004basic}
means to attain dependability 
Fault Prevention
- prevention techniques like strong typing, information hiding, modularity, 
Fault Tolerance
- error detection and system reovery, corrective maintenance techniques
Error Detection -  concurrent detection (takes place during service) and preemptive detection (takes place while operation is suspended)
- Recovery - error handling and fault handling
- Error Handling (eliminates errors) - rollback (go back to saved state before error), rollforward (go to state without the error), compensation (erroneous state contains enough info the be masked, conceals a progeessive loss)
- fault handling (prevents faults from occuring again) - diagnosis (indentify the cause and location), isolation (exclude the faulty component), reconfiguration (swap in spare components or reassign task), reinitialization (updates and records the new configuration)
Fault Removal 
- consist of verfication before and and after removal to ensure that removal does not cause any consequences (Static verfication =  static analysis or theorem proving, or model checking) (Dynamic verfication = testing through symbol or real inputs)
- fault removal during use - corrective or preventative maintainance
Fault Forecasting
- evlauation of system behaviour with respect to fault occurance and activation in a qualitative and quantitative method


\cite{benveniste2012contracts}
devised a mathematical meta theory of contracts, with compositions of components formed through contracts
components are abstract entities that can be compused through associative or commutative links
Contracts define a set of valid implementations and environments it can work in to contrain both sides of the interaction
Contracts have assumptions and guaruntees
observers can be expressed through temporal logic, simulink, modelica for runtime monitering
contract meta theory defines autonomous components, with composability (typed) and integration through contracts

Contract-based design, as developed by \citet{benveniste2012contracts}, treats composition as a mechanism for safely decomposing systems into subsystems under strong assumptions about interface and envirnonmental stability. In these settings contracts are used to ensure that independently developed components can be recomposed without violating correctness.

In contrast, Systems of Twinned Systems operate under fundamentally different conditions. Constituents are loosely coupled, interact through asynchronous event streams, and evolve independently, leading to environments in which contract assumptions are routinely violated. Under contract based design, such systems would be deemed incompatible or inconsistent, and meaningful computation would not be possible.

This thesis addresses the inverse problem: rather than using contracts to prevent unreliable composition, it investigates how computation can proceed despite it. By shifting reliability concerns from the communication layer to the computational semantics of event processing, and by explicitly representing uncertainty through confidence-aware reconstruction and pattern detection, the proposed architecture enables meaningful reasoning in dynamically composed, unreliable systems.